% =====================================================================
\chapter{Entropía relativa e información mutua}
\label{tema:2}
% =====================================================================

Este tema introduce la entropía relativa o divergencia de Kullback--Leibler, que mide el coste de usar una distribución equivocada, y la información mutua, que mide cuánto informa una variable sobre otra. El resultado que lo articula todo es la desigualdad de Gibbs, que se deduce de la desigualdad de Jensen. Se sigue~\cite[§2.3--2.6 y §2.8]{ct}; la interpretación de las dos direcciones de la divergencia en aprendizaje automático procede del tema «Comparing probability distributions and machine learning» de~\cite{camit}, y la de la información mutua, de~\cite[cap.~8]{mackay}.

\section{La desigualdad de Jensen}

\begin{definicion}[Función convexa]\label{def:convexa}
Sea $I\subseteq\R$ un intervalo. Una función $f\colon I\to\R$ es \emph{convexa} si
\[
f\bigl(\lambda a+(1-\lambda)b\bigr)\le\lambda f(a)+(1-\lambda)f(b)\qquad\text{para todo }a,b\in I,\ \lambda\in[0,1],
\]
y \emph{estrictamente convexa} si la desigualdad es estricta cuando $a\ne b$ y $\lambda\in(0,1)$. La función $f$ es (estrictamente) \emph{cóncava} si $-f$ es (estrictamente) convexa.
\end{definicion}

Si $f$ es dos veces derivable con $f''\ge 0$ ($f''>0$) en $I$, entonces es convexa (estrictamente convexa)~\cite[§2.6]{ct}; el módulo~4 estudia este criterio. Así, $\ln$ y $\log$ son estrictamente cóncavas en $(0,\infty)$, porque $(\ln t)''=-1/t^{2}<0$, y $t\mapsto t\log t$ es estrictamente convexa en $(0,\infty)$.

\begin{teorema}[Desigualdad de Jensen]\label{teo:jensen}
Sea $f$ convexa en un intervalo $I$ y sea $X$ una variable aleatoria con valores en $I$ que toma un número finito de valores. Entonces
\[
\E[f(X)]\ge f\bigl(\E[X]\bigr).
\]
Si $f$ es estrictamente convexa, la igualdad se da si y solo si $X$ es constante con probabilidad~1. Para $f$ cóncava la desigualdad se invierte.
\end{teorema}

\section{Entropía relativa}

\begin{definicion}[Entropía relativa]\label{def:kl}
Sean $p$ y $q$ distribuciones sobre $\X$. La \emph{entropía relativa} o \emph{divergencia de Kullback--Leibler} de $p$ respecto de $q$ es
\[
\KL{p}{q}=\sum_{x\in\X}p(x)\log\frac{p(x)}{q(x)}=\E_{X\sim p}\Bigl[\log\frac{p(X)}{q(X)}\Bigr],
\]
con los convenios $0\log\frac{0}{q}=0$ y $p\log\frac{p}{0}=+\infty$ si $p>0$.
\end{definicion}

En particular, $\KL{p}{q}=+\infty$ en cuanto $q$ declara imposible un resultado que, según $p$, puede ocurrir. Es la versión informacional de la pérdida infinita del módulo~1 cuando un modelo asigna probabilidad~0 a un token del texto.

\begin{teorema}[Desigualdad de Gibbs o de la información]\label{teo:gibbs}
Para distribuciones $p$ y $q$ sobre $\X$, $\KL{p}{q}\ge 0$, con igualdad si y solo si $p=q$.
\end{teorema}

\addtocounter{definicion}{1}%

\begin{observacion}
La entropía relativa no es simétrica y no cumple la desigualdad triangular; por eso se llama \emph{divergencia} y no distancia. Sí cumple la propiedad que importa para comparar un modelo con los datos: es no negativa y solo se anula cuando las dos distribuciones coinciden.
\end{observacion}

\section{Entropía cruzada}

\begin{definicion}[Entropía cruzada]\label{def:cruzada}
La \emph{entropía cruzada} de $q$ respecto de $p$ es $H(p,q)=-\sum_{x}p(x)\log q(x)=\E_{X\sim p}[-\log q(X)]$.
\end{definicion}

\begin{proposicion}[Descomposición de la entropía cruzada]\label{prop:descomposicion}
$H(p,q)=H(p)+\KL{p}{q}$. En consecuencia, $H(p,q)\ge H(p)$, con igualdad si y solo si $q=p$.
\end{proposicion}

En el \cref{ej:asimetria}, $H(p)=1{,}5$ bits y $H(p,q)=1{,}822$ bits, cuya diferencia es $\KL{p}{q}=0{,}322$ bits. El \cref{tema:3} dará a esta identidad un sentido operativo: $H(p,q)$ es la longitud media de un código diseñado para $q$ cuando los datos siguen $p$, y $\KL{p}{q}$ son los bits que se desperdician por usar el modelo equivocado.

\addtocounter{definicion}{1}%

\begin{observacion}[Las dos direcciones de la divergencia]\label{obs:direcciones}
Al ajustar una distribución $q$ a una distribución $p$ dada pueden minimizarse dos cantidades distintas. La divergencia \emph{hacia delante} $\KL{p}{q}$ penaliza con fuerza que $q(x)$ sea pequeña donde $p(x)$ no lo es: $q$ tiende a \emph{cubrir} todo el soporte de $p$. Es la que se minimiza al entrenar por máxima verosimilitud. La divergencia \emph{hacia atrás} $\KL{q}{p}$ penaliza que $q$ ponga masa donde $p$ es pequeña: $q$ tiende a \emph{concentrarse} en una de las modas de $p$. Es la que aparece en la inferencia variacional y en la penalización del ajuste por preferencias del \cref{tema:5}~\cite{camit}.
\end{observacion}

\section{Información mutua}

\begin{definicion}[Información mutua]\label{def:mutua}
La \emph{información mutua} entre $X$ e $Y$ es la entropía relativa entre la distribución conjunta y el producto de las marginales:
\[
I(X;Y)=\KL{p(x,y)}{p(x)\,p(y)}=\sum_{x,y}p(x,y)\log\frac{p(x,y)}{p(x)\,p(y)}.
\]
\end{definicion}

\begin{proposicion}\label{prop:mutua-identidades}
\[
I(X;Y)=H(X)+H(Y)-H(X,Y)=H(Y)-H(Y\mid X)=H(X)-H(X\mid Y).
\]
En particular, $I(X;Y)=I(Y;X)$ e $I(X;X)=H(X)$.
\end{proposicion}

La demostración de la primera igualdad se propone en el \cref{eje:2-mutua}; las otras dos se siguen de ella y de la regla de la cadena (\cref{teo:cadena-entropia}). Para $Y=X$ se tiene $H(X\mid X)=0$.

\begin{teorema}\label{teo:mutua-positiva}
$I(X;Y)\ge 0$, con igualdad si y solo si $X$ e $Y$ son independientes. En consecuencia, $H(Y\mid X)\le H(Y)$: condicionar no aumenta, en media, la entropía.
\end{teorema}

\begin{observacion}
La desigualdad $H(Y\mid X)\le H(Y)$ vale en media, no para cada valor: puede ocurrir $H(Y\mid X=x)>H(Y)$ para algún $x$ concreto. Una observación sorprendente puede aumentar la incertidumbre, aunque observar $X$ la reduce en promedio~\cite[§2.6]{ct}.
\end{observacion}

\addtocounter{figure}{1}%

\addtocounter{definicion}{1}%

\addtocounter{definicion}{1}%

\section{La desigualdad del procesamiento de datos}

\begin{definicion}[Información mutua condicionada]
Para variables $X$, $Y$ y $Z$, la \emph{información mutua entre $X$ e $Y$ condicionada a $Z$} es $I(X;Y\mid Z)=H(X\mid Z)-H(X\mid Y,Z)$, donde $H(X\mid Y,Z)$ es la entropía de $X$ condicionada al par $(Y,Z)$.
\end{definicion}

Se cumple $I(X;Y\mid Z)=\sum_z p(z)\,I(X;Y\mid Z=z)\ge 0$, porque cada sumando es la información mutua de la distribución condicionada a $Z=z$ y es no negativo por el \cref{teo:mutua-positiva}. Además, aplicando la regla de la cadena a $H(X)-H(X\mid Y,Z)=\bigl[H(X)-H(X\mid Z)\bigr]+\bigl[H(X\mid Z)-H(X\mid Y,Z)\bigr]$ se obtiene la \emph{regla de la cadena para la información mutua}~\cite[§2.5]{ct}:
\[
I(X;Y,Z)=I(X;Z)+I(X;Y\mid Z).
\]

\begin{definicion}[Cadena de Markov]
Las variables $X$, $Y$ y $Z$ forman una \emph{cadena de Markov} $X\to Y\to Z$ si $p(z\mid x,y)=p(z\mid y)$ para todo $x,y,z$; es decir, si $Z$ depende de $X$ solo a través de $Y$. Esto equivale a que $X$ y $Z$ sean condicionalmente independientes dado $Y$, y en ese caso $I(X;Z\mid Y)=0$.
\end{definicion}

\begin{teorema}[Desigualdad del procesamiento de datos]\label{teo:procesamiento}
Si $X\to Y\to Z$, entonces $I(X;Z)\le I(X;Y)$.
\end{teorema}

\begin{corolario}\label{cor:procesamiento-funcion}
Para toda función $g$, $I(X;g(Y))\le I(X;Y)$.
\end{corolario}

\begin{observacion}
Ningún procesamiento de un dato, por elaborado que sea, aumenta la información que contiene sobre otra variable. Las representaciones internas de una red neuronal son funciones de su entrada; por el \cref{cor:procesamiento-funcion}, no contienen más información sobre la respuesta correcta que la propia entrada. Lo que el entrenamiento consigue es hacer esa información \emph{accesible} para una capa de salida sencilla, no crearla.
\end{observacion}

